% Generated by scripts/materialize_unified_paper.py; do not edit.
\section{Prior art and comparison with a reported lower bound}
\label{qua:sec:prior-art}

Previous quantum papers under the phrase ``Ramsey algorithm'' solve different
problems.  Gaitan and Clark encode the existence of avoidance colourings in
an adiabatic Hamiltonian \cite{GaitanClark2012}, and Bian et al.\ implement
small instances experimentally \cite{BianEtAl2013}; Wang searches the space of
$2^{\binom V2}$ labelled graphs when testing a Ramsey threshold
\cite{Wang2016}; Qu et al.\ use quantum counting for hypergraph Ramsey numbers
\cite{QuLiWangBaoCao2013}; and generalized/ordinary bound discovery has been
studied through adiabatic optimization and QUBO reduction
\cite{RanjbarMacreadyClarkGaitan2016,PionMniszewski2025}.  Their square roots
refer to the number of graph colourings, not the number of vertices of one
given graph.  D{\"o}rn's clique and independent-set query algorithms are
closer in oracle model but do not exploit the Ramsey totality promise or the
majority recursion \cite{Doern2005}.  Quantum subset-finding algorithms already
give nontrivial clique-query bounds \cite{ChildsEisenberg2005}, while
learning-graph algorithms for fixed subgraph containment
\cite{LeeMagniezSantha2012,Zhu2012} and quantum backtracking
\cite{Montanaro2018}
provide powerful generic primitives.  None by itself supplies the
implicit-set survival invariant used here.  To our knowledge, no earlier work uses
implicit-set quantum sampling to implement the constructive
Erdős--Szekeres proof; this priority statement is limited to the exact oracle
model and sources compared here.

Kraj{\'\i}\v{c}ek connected Ramsey and weak-pigeonhole proof complexity
\cite{Krajicek2001} and later formulated the structured-pigeonhole search
problem for circuit-encoded homogeneous sets \cite{Krajicek2005}.  Pasarkar,
Papadimitriou, and Yannakakis
place Ramsey search in their polynomial long-choice class PLC, which captures
iterated pigeonhole arguments \cite{PasarkarPapadimitriouYannakakis2023}.
\paragraph{Comparison at the formal default parameters.}
The closest exact-model source is the FOCS 2024 proceedings paper of
Jain--Li--Robere--Xun \cite{JainLiRobereXun2024}.  Definition II.6 of its
proceedings version (p.~413) uses $N=2^n$ vertices, and the convention after
Definition II.7 sets the default target to $K=n/2$.  We use these formal
definitions throughout this comparison.  The unnumbered ``Query
complexity of RAMSEY'' paragraph in Section III (p.~415) states an $N^{1-o(1)}$
quantum lower bound via Lemma III.1 and the Liu--Zhandry multicollision theorem
\cite{LiuZhandry2019}.  The graph-hash product in Lemma III.1 can be written
off the diagonal as
$A_h(u,v)=\mathbf1[h(u)=h(v)]\lor A_0(h(u),h(v))$, with
$A_h(u,u)=0$.  It is a valid symmetric graph and has an $O(1)$-query coherent
implementation from the value oracle for $h$, using computation and
uncomputation of its two values.  Thus the promise of a valid graph does not
exclude this reduction.  Corollary~\ref{qua:cor:unpromised} also aligns our upper bound with the
locally verifiable arbitrary-circuit relation.  Moreover,
\cref{qua:cor:succinct} returns $\lfloor n/2\rfloor+1$ vertices, so discarding one
if necessary solves their default target up to the suppressed rounding
convention.  Thus malformed encodings and the target convention do not remove
the apparent conflict.

\paragraph{Direct parameter substitution.}
Write $H$ for the multicollision range size and $G$ for the number of
vertices in the Ramsey instance.  For fixed multiplicity $t$, the sufficient
parameter condition in Jain et al.'s Theorem I.3 is
$G\ge H^{4t}/4^t$.  At the boundary
$G=\Theta_t(H^{4t})$, Liu--Zhandry's exponent in $H$ is
$\alpha_t=(2^{t-1}-1)/(2^t-1)$.  The inherited exponent is therefore
\[
 \frac{\alpha_t}{4t}
 =\frac{2^{t-1}-1}{4t(2^t-1)}\le\frac1{24},
 \qquad t\ge2,
\]
with equality at $t=2$.  Indeed, $t=2$ gives $1/24$, while for $t\ge3$
the strict inequality $\alpha_t<1/2$ gives
$\alpha_t/(4t)<1/(8t)\le1/24$.  Increasing $G$ relative to $H$ cannot
improve this transferred exponent.  The cited Liu--Zhandry result is for
fixed $t$ and does not by itself permit substitution of a growing
multiplicity.  This calculation concerns what follows from these stated
ingredients; it is not an upper bound on the true quantum query complexity.

Consequently, the stated near-linear consequence is not obtained by directly
composing the two cited statements.  If that consequence is intended for the
same bounded-error quantum-query relation, it is incompatible with our upper
bound.  Our proof does not use this consequence or the Liu--Zhandry lower
bound.  This observation concerns
only the parameter substitution in the unnumbered query-complexity discussion
and does not concern the principal black-box nonreducibility or TFNP separation
theorems of Jain et al.  An additional ingredient or a different intended
normalization may resolve the discrepancy.
